\chapter{Theses \& Dissertations}
\label{ch:theses}

This chapter covers four \texttt{scrbook}-based doctypes for long-form
academic works: \texttt{thesis}, \texttt{dissertation},
\texttt{thesis-tuhh}, and \texttt{dictionary}.  The \texttt{manual} doctype
uses \texttt{scrreprt} and is documented in
Chapter~\ref{ch:reports}.

All book-type doctypes inherit \texttt{listof=totoc},
\texttt{bibliography=totoc}, \texttt{chapterprefix=true},
\texttt{open=right}, and \texttt{numbers=noenddot}.  They support
\texttt{\textbackslash frontmatter}, \texttt{\textbackslash mainmatter},
and \texttt{\textbackslash backmatter} for distinct typographic treatment
of preliminary pages, body chapters, and appendices.

\section{Thesis vs.\ Dissertation}

\begin{table}[htbp]
    \centering
    \caption{Comparison of thesis and dissertation doctypes}
    \label{tab:thesis-vs-dissertation}
    \begin{tabular}{@{} l p{5.5cm} p{5.5cm} @{}}
        \toprule
        & \textbf{thesis} & \textbf{dissertation} \\
        \midrule
        \textbf{KOMA base}   & \texttt{scrbook} & \texttt{scrbook} \\[3pt]
        \textbf{Font size}   & 12\,pt & Default (11\,pt) \\[3pt]
        \textbf{Sides}       & Two-sided & Two-sided \\[3pt]
        \textbf{secnumdepth} & 2 & 3 \\[3pt]
        \textbf{Line spacing} & One-and-a-half & Default \\[3pt]
        \textbf{Paragraph spacing} & None & Default \\[3pt]
        \textbf{Supervisor model}  & Single advisor & Primary + secondary \\[3pt]
        \textbf{Committee}   & \texttt{\textbackslash thesiscommittee} &
                               \texttt{\textbackslash dissertationcommittee} \\[3pt]
        \textbf{Degree cmd}  & \texttt{\textbackslash thesisdegree} &
                               \texttt{\textbackslash dissertationdegree} \\[3pt]
        \textbf{Keywords}    & No & Yes (\texttt{\textbackslash keywords}) \\
        \bottomrule
    \end{tabular}
\end{table}

The \texttt{thesis} doctype is designed for master's theses with a single
advisor, shallower sectioning depth, and larger 12\,pt font.  The
\texttt{dissertation} doctype targets doctoral works with deeper sectioning,
dual supervisor support, and keyword fields for library cataloguing.

\section{Doctype Profiles}

\subsection{thesis}

The \texttt{thesis} doctype produces a formal thesis title page with
institution, department, degree, advisor, committee, and defense date.  It
uses 12\,pt serif font, one-and-a-half line spacing, and no paragraph
indentation.  The document is two-sided with binding correction support via
the \texttt{BCOR} option.

\begin{table}[htbp]
    \centering
    \caption{Custom commands for \texttt{thesis}}
    \label{tab:thesis-commands}
    \begin{tabular}{@{} l p{8cm} @{}}
        \toprule
        \textbf{Command} & \textbf{Description} \\
        \midrule
        \texttt{\textbackslash thesisinstitution\{...\}}  & University or institution name \\
        \texttt{\textbackslash thesisdepartment\{...\}}   & Department or school name \\
        \texttt{\textbackslash thesisdegree\{...\}}       & Degree (default: Master of Science) \\
        \texttt{\textbackslash thesisadvisor\{...\}}      & Thesis advisor name \\
        \texttt{\textbackslash thesiscommittee\{...\}}    & Committee members (comma-separated) \\
        \texttt{\textbackslash defensedate\{...\}}        & Defense date (default: \texttt{\textbackslash today}) \\
        \texttt{\textbackslash thesislocation\{...\}}     & City and country \\
        \bottomrule
    \end{tabular}
\end{table}

Use \texttt{thesis} for master's theses or similar graduate-level documents
that require a formal academic title page.

\begin{minted}{latex}
\documentclass[doctype=thesis, language=english]{../../omnilatex}
\title{Numerical Methods for Turbulent Flow Simulation}
\author{Jane Doe}
\thesisinstitution{Technical University}
\thesisdepartment{Mechanical Engineering}
\thesisdegree{Master of Science}
\thesisadvisor{Prof.\ Dr.\ A.~Mueller}
\thesiscommittee{Prof.\ B.~Smith, Prof.\ C.~Lee}
\defensedate{2025-07-15}
\thesislocation{Hamburg, Germany}
\begin{document}
\maketitle
\frontmatter
\tableofcontents
\mainmatter
\chapter{Introduction}
\end{document}
\end{minted}

\subsection{dissertation}

The \texttt{dissertation} doctype extends the thesis format with a dual
supervisor model (primary and secondary), deeper sectioning
(\texttt{secnumdepth=3}), and keyword support for repository metadata.  The
title page includes institution, department, degree, both supervisors,
committee, submission date, and keywords.

\begin{table}[htbp]
    \centering
    \caption{Custom commands for \texttt{dissertation}}
    \label{tab:dissertation-commands}
    \begin{tabular}{@{} l p{8cm} @{}}
        \toprule
        \textbf{Command} & \textbf{Description} \\
        \midrule
        \texttt{\textbackslash dissertationdegree\{...\}}        & Degree (default: Doctoral Dissertation) \\
        \texttt{\textbackslash dissertationdepartment\{...\}}     & Department or faculty \\
        \texttt{\textbackslash dissertationinstitution\{...\}}   & University name \\
        \texttt{\textbackslash firstsupervisor\{...\}}           & Primary supervisor \\
        \texttt{\textbackslash secondsupervisor\{...\}}          & Secondary or co-supervisor \\
        \texttt{\textbackslash dissertationcommittee\{...\}}     & Advisory committee members \\
        \texttt{\textbackslash submissiondate\{...\}}            & Submission date (default: \texttt{\textbackslash today}) \\
        \texttt{\textbackslash keywords\{...\}}                  & Comma-separated keywords \\
        \bottomrule
    \end{tabular}
\end{table}

Use \texttt{dissertation} for PhD or doctoral works that require dual
supervisor tracking and keyword metadata.

\begin{minted}{latex}
\documentclass[doctype=dissertation]{../../omnilatex}
\title{Data-Driven Modelling of Urban Microclimates}
\author{Jane Doe}
\dissertationinstitution{Technical University}
\dissertationdepartment{Environmental Engineering}
\dissertationdegree{Doctor of Engineering}
\firstsupervisor{Prof.\ Dr.\ A.~Mueller}
\secondsupervisor{Prof.\ Dr.\ B.~Smith}
\dissertationcommittee{Dr.\ C.~Lee, Dr.\ D.~Park}
\submissiondate{2025-06-01}
\keywords{urban microclimate, data-driven modelling, CFD}
\begin{document}
\maketitle
\frontmatter
\chapter{Abstract}
\tableofcontents
\mainmatter
\chapter{Introduction}
\end{document}
\end{minted}

\subsection{thesis-tuhh}

The \texttt{thesis-tuhh} doctype combines the generic \texttt{thesis} profile
with the \texttt{institution=tuhh} option for institution-specific branding
at Hamburg University of Technology (TUHH).  It is not a separate profile
file; rather, it is activated by combining options:

\begin{minted}{latex}
\documentclass[
    doctype=thesis,
    institution=tuhh,
    titlestyle=TUHH,
    loadGlossaries,
    todonotes,
]{../../omnilatex}
\end{minted}

Key features of the TUHH configuration include:

\begin{itemize}
    \item \textbf{Title page branding} -- The \texttt{titlestyle=TUHH}
          option selects a custom title page layout with the TUHH logo,
          bilingual headers, and the institutional colour scheme.
    \item \textbf{Bilingual logos} -- Both German and English TUHH logos
          are available, selected automatically based on the
          \texttt{language} option.
    \item \textbf{TUHHuman TikZ picture} -- A TikZ-drawn mascot
          (\texttt{ TUHHuman}) is defined in the TUHH institution package
          and can be placed on the title page or in colophon sections.
    \item \textbf{Custom title page} -- The TUHH institution package
          (\texttt{config/institutions/tuhh.sty}) overrides the default
          \texttt{\textbackslash maketitle} with a layout that includes the
          TUHH wordmark, faculty name, and thesis metadata in the
          institutional style.
\end{itemize}

The complete TUHH thesis example is located at
\texttt{examples/thesis-tuhh/} and demonstrates front matter (abstract,
task description, authorship declaration), main matter chapters with code
listings and figures, and back matter with appendices.

Use \texttt{thesis-tuhh} for any thesis produced at TUHH that requires
official institutional branding on the title page.

\begin{minted}{latex}
\documentclass[
    language=english,
    institution=tuhh,
    titlestyle=TUHH,
    censoring=true,
    loadGlossaries,
    todonotes,
]{../../omnilatex}
\RequirePackage{config/document-settings}
\addbibresource{bib/bibliography.bib}
\begin{document}
\import{content/}{frontmatter}
\import{content/}{mainmatter}
\import{content/}{backmatter}
\end{document}
\end{minted}

\subsection{dictionary}

The \texttt{dictionary} doctype produces a reference dictionary or glossary
book.  It uses a compact 10\,pt font with \texttt{DIV=13} for a dense,
reference-style layout.  The title page includes series name, subtitle,
edition, publisher, ISBN, and subject fields.  The document is two-sided
with shallow sectioning (\texttt{secnumdepth=1}, \texttt{tocdepth=1}).

\begin{table}[htbp]
    \centering
    \caption{Custom commands for \texttt{dictionary}}
    \label{tab:dictionary-commands}
    \begin{tabular}{@{} l p{8cm} @{}}
        \toprule
        \textbf{Command} & \textbf{Description} \\
        \midrule
        \texttt{\textbackslash dictionaryseries\{...\}}     & Series collection name \\
        \texttt{\textbackslash dictionarysubtitle\{...\}}   & Subtitle or tagline \\
        \texttt{\textbackslash dictionaryedition\{...\}}    & Edition (default: First Edition) \\
        \texttt{\textbackslash dictionarypublisher\{...\}}  & Publisher name \\
        \texttt{\textbackslash dictionaryisbn\{...\}}       & ISBN \\
        \texttt{\textbackslash dictionarysubjects\{...\}}  & Subject categories \\
        \texttt{\textbackslash dictionarynote\{...\}}       & Editorial note or disclaimer \\
        \bottomrule
    \end{tabular}
\end{table}

Use \texttt{dictionary} for compiling terminology collections, technical
lexicons, or domain-specific reference books.

\begin{minted}{latex}
\documentclass[doctype=dictionary]{../../omnilatex}
\title{Dictionary of Control Engineering}
\author{Jane Doe}
\dictionaryseries{OmniLexicon Collection}
\dictionarysubtitle{A quick reference for control systems terminology}
\dictionaryedition{Second Edition}
\dictionarypublisher{Acme Academic Press}
\dictionaryisbn{978-0-00-000002-4}
\dictionarysubjects{Control Theory, Automation}
\begin{document}
\maketitle
\chapter{A}
\section{Actuator}
A device that converts a control signal into mechanical motion.
\section{Adaptive Control}
A control strategy that adjusts parameters in real time.
\end{document}
\end{minted}

\subsection{manual (book-class)}

The \texttt{manual} doctype uses \texttt{scrreprt} and is documented in
Chapter~\ref{ch:reports}.  When used as part of a book-style workflow with
\texttt{\textbackslash frontmatter} and \texttt{\textbackslash backmatter},
it integrates naturally with the book-class ecosystem.  See
\Cref{tab:manual-commands} for its custom commands.
